%--------------------------------------------------------- general introduction
\cleardoublepage
\phantomsection
\addcontentsline{toc}{chapter}{General Introduction}
\markboth{General Introduction}{}
\chapter*{General Introduction}

\section*{Context}

Online retail has made large catalogues cheap to publish and difficult to
navigate. A physical shop solves this with a person: a customer describes a
problem in their own words, and an assistant translates that description into
a small set of suitable products. An online catalogue of a thousand items
offers categories and a search box instead, and leaves the translation entirely
to the customer.

The gap is particularly wide in cosmetics and natural health products, where
suitability depends on attributes the customer may not know how to name.
A shopper who says \emph{``I have dark spots since my pregnancy and my skin
reacts to everything''} is describing hyperpigmentation on sensitive skin. No
category menu captures that sentence, and no keyword search resolves it into a
coherent routine of cleanser, treatment and protection.

This project was carried out at \companyname, a company that develops and
ships WordPress applications. The work addressed a live client catalogue: a
Tunisian online shop selling natural cosmetics, dietary supplements and organic
food, running on WooCommerce with more than six hundred published products.

\section*{Problem statement}

An initial audit of the client database revealed that the information required
to make recommendations was almost entirely absent from the structured data:

\begin{itemize}
  \item the catalogue contained \textbf{zero} product tags;
  \item the skin-type attribute was populated on \textbf{six products out of
        six hundred and twelve};
  \item product categories described physical format (face, hair, body) rather
        than customer need, and included merchandising buckets such as
        \emph{Homepage} holding 239 products;
  \item the packs sold as ready-made routines were stored as ordinary products,
        with no database link to their contents.
\end{itemize}

The information did exist, but only as free text: product descriptions
averaging 1{,}414 characters, written to a consistent commercial pattern and
routinely stating the skin types and concerns a product addresses.

The problem therefore divided into two: extracting structure from unstructured
commercial text, and building a recommendation layer that a non-technical
shopper can actually use.

\section*{Objectives}

The project set out to:

\begin{enumerate}
  \item audit the production catalogue and establish what recommendation-grade
        data was genuinely available;
  \item design a product data model suited to recommendation rather than to
        catalogue display;
  \item build an automated pipeline extracting that model from the live
        database, repeatable at no marginal cost;
  \item implement a deterministic, explainable recommendation engine producing
        ordered skincare routines;
  \item expose the engine as a hosted service with live pricing and stock;
  \item add a natural-language layer allowing a shopper to describe a need in
        ordinary French;
  \item extend that layer into a conversational agent capable of multi-turn
        reasoning over the catalogue;
  \item validate the recommendations against realistic customer profiles.
\end{enumerate}

\section*{Report outline}

\textbf{Chapter 1 --- General Context and Environment of the Project} presents the host enterprise (\companyname), details the e-commerce landscape in Tunisia, outlines the catalogue data constraints, formalises the core engineering problem, and details the project methodology and temporal timeline.

\textbf{Chapter 2 --- Market Study and State of the Art} examines commercial personalisation in the beauty and wellness sector, reviews the theoretical and technical foundations of Natural Language Processing, information extraction, Large Language Models, and ReAct tool use, contrasts existing solutions, defines the proposed hybrid architecture, and details the adopted CRISP-DM methodology.

\textbf{Chapter 3 --- Methodology, Data Collection, and Analytical Framework} details the end-to-end data transformation workflow, presents the exploratory auditing and entity-relationship schema of the WooCommerce database, quantifies textual completeness and category distributions, describes the extraction and TF-IDF pack recovery preprocessing, formalises the derived multi-dimensional product data model, and presents the deterministic scoring mathematics.

\textbf{Chapter 4 --- Implementation, Deployment, and Agent Evaluation} describes the implementation of the system components (extraction pipeline, scoring engine, asynchronous FastAPI microservice, WordPress custom shortcode plugin, and interactive web demo interfaces), presents the Docker Compose containerisation, and establishes a comprehensive empirical evaluation protocol over 22 realistic customer scenarios, analysing prompt sensitivity, performance metrics, failure modes, and industrial positioning.

The \textbf{General Conclusion and Perspectives} synthesises the principal contributions, examines the engineering challenges encountered, and outlines promising future directions for production deployment and research.

\cleardoublepage
